\chapter{Complexity — ND Clock Host Run}
%%% generated by the blueprint pipeline from TCSlib.Complexity.TuringMachine.Build.NDClockHostRun
\section{Overview}

Declarations of the Lean module \texttt{TCSlib.Complexity.TuringMachine.Build.NDClockHostRun}.

\section{Declarations}

\begin{theorem}[Splitting a Boolean Word]
\lean{Turing.ClockHost.exists_split_block}
\label{Turing.ClockHost.exists_split_block}
\leanok
\difficulty{4}
If a Boolean word $w$ has length $m+1+r$ for natural numbers $m$ and $r$, then there are
a Boolean word $u$, a Boolean $b$, and a Boolean word $v$ such that $w=u \, ++ \, [b] \,
++ \, v$, with $|u|=m$ and $|v|=r$.
\end{theorem}

\begin{theorem}[Recursive Leftward Head Movement]
\lean{Turing.ClockHost.inLeft_succ}
\label{Turing.ClockHost.inLeft_succ}
\leanok
\uses{Turing.ClockHost.inLeft, Turing.moveInputPos}
\difficulty{4}
For every natural number $m$, every $n$, every input-head position $ip$ among the $n+2$
positions, and every natural number $j$, the result of $j+1$ successive leftward moves
from $ip$ equals one leftward move applied to the result of $j$ successive leftward
moves from $ip$, with moves clamped to the valid positions.
\end{theorem}

\begin{theorem}[Leftward clock sweep on a marked cell]
\lean{Turing.ClockHost.p0_mark}
\label{Turing.ClockHost.p0_mark}
\leanok
\uses{Turing.ClockHost.apply_hostCfg, Turing.ClockHost.clockTr, Turing.ClockHost.hostCfg, Turing.ClockHost.hostCfg_symbols_yI, Turing.ClockHost.hostTM, Turing.ClockHost.stepWith_hostCfg, Turing.ClockHost.yard, Turing.NDTM.stepWith, Turing.moveInputPos}
\difficulty{3}
Let $S$ be a deterministic machine, $A$ a nondeterministic machine, and $x$ a Boolean
input word; let $m\in\mathbb N$, let $b$ be a choice bit, and choose an input-head
position, work-tape contents and head positions for both machines, and an integer $z$.
If the yard tape contains true at cell $z$, then one host-machine step in phase $p0$
moves the input head left with clamping, advances each of $A$’s work-tape heads by one,
moves the yard-tape head from $z$ to $z-1$, and leaves the phase, all tape contents, the
$S$ work-tape head positions, and the output unchanged.
\end{theorem}

\begin{theorem}[Turn from the first to second reset sweep]
\lean{Turing.ClockHost.p0_turn}
\label{Turing.ClockHost.p0_turn}
\leanok
\uses{Turing.ClockHost.apply_hostCfg, Turing.ClockHost.clockTr, Turing.ClockHost.hostCfg, Turing.ClockHost.hostCfg_symbols_yI, Turing.ClockHost.hostTM, Turing.ClockHost.stepWith_hostCfg, Turing.ClockHost.yard, Turing.ClockHost.yard_neg_one, Turing.NDTM.stepWith, Turing.moveInputPos}
\difficulty{3}
Let $S$ be a deterministic machine, $A$ a nondeterministic machine, and $x$ a Boolean
input word; let $m\in\mathbb N$, let $b$ be a choice bit, and choose an input-head
position and work-tape contents and head positions for both machines. One host-machine
step from phase $p0$ with yard-tape head at $-1$ moves the input head left with
clamping, advances each of $A$’s work-tape heads by one, moves the yard-tape head to
$0$, and changes the phase to $p1$, leaving all tape contents, the $S$ work-tape head
positions, and the output unchanged.
\end{theorem}

\begin{theorem}[Rightward sweep mark step]
\lean{Turing.ClockHost.p1_mark}
\label{Turing.ClockHost.p1_mark}
\leanok
\uses{Turing.ClockHost.apply_hostCfg, Turing.ClockHost.clockTr, Turing.ClockHost.hostCfg, Turing.ClockHost.hostCfg_symbols_yI, Turing.ClockHost.hostTM, Turing.ClockHost.stepWith_hostCfg, Turing.ClockHost.yard, Turing.NDTM.stepWith}
\difficulty{3}
Let $S$ be a deterministic Boolean-tape Turing machine and $A$ a nondeterministic
Boolean-tape Turing machine, with a Boolean loop parameter and input word $x$. For any
natural number $m$, choice bit, input-head position, tape contents and head positions,
and integer $z$, if the yard tape of length $m$ contains a mark at $z$, then one
host-machine step from phase $p1$ leaves the phase and input head unchanged, advances
every head of $A$’s work tapes by one, advances the yard head by one, and leaves the
other tapes and output unchanged.
\end{theorem}

\begin{theorem}[Turnaround of the second reset sweep]
\lean{Turing.ClockHost.p1_turn}
\label{Turing.ClockHost.p1_turn}
\leanok
\uses{Turing.ClockHost.apply_hostCfg, Turing.ClockHost.clockTr, Turing.ClockHost.hostCfg, Turing.ClockHost.hostCfg_symbols_yI, Turing.ClockHost.hostTM, Turing.ClockHost.stepWith_hostCfg, Turing.ClockHost.yard, Turing.ClockHost.yard_len, Turing.NDTM.stepWith}
\difficulty{3}
Let $S$ be a deterministic Turing machine over the Boolean alphabet, let $A$ be a
nondeterministic Turing machine over the Boolean alphabet, let $\ell$ be a Boolean, and
let $x$ be a finite Boolean input word. For every $m\in\mathbb{N}$, choice bit $b$,
valid input-head position $ip$, work-tape contents and head positions for both machines,
one step of the host machine from phase $p1$, with the yard tape containing $m$
consecutive true symbols starting at cell $0$ and its head at cell $m$, leaves the
phase-specific tapes of $S$ and the output unchanged, advances each work-tape head of
$A$ by one, leaves the input-head position unchanged, keeps the yard tape unchanged, and
enters phase $p2$ with the yard head at cell $m-1$.
\end{theorem}

\begin{theorem}[Leftward sweep mark step]
\lean{Turing.ClockHost.p2_mark}
\label{Turing.ClockHost.p2_mark}
\leanok
\uses{Turing.ClockHost.apply_hostCfg, Turing.ClockHost.clockTr, Turing.ClockHost.hostCfg, Turing.ClockHost.hostCfg_symbols_yI, Turing.ClockHost.hostTM, Turing.ClockHost.stepWith_hostCfg, Turing.ClockHost.yard, Turing.NDTM.stepWith}
\difficulty{3}
Let $S$ be a deterministic Boolean-tape Turing machine and $A$ a nondeterministic
Boolean-tape Turing machine, with a Boolean loop parameter and input word $x$. For any
natural number $m$, choice bit, input-head position, tape contents and head positions,
and integer $z$, if the yard tape of length $m$ contains a mark at $z$, then one
host-machine step from phase $p2$ leaves the phase and input head unchanged, advances
every head of $A$’s work tapes by one, moves the yard head to $z-1$, and leaves the
other tapes and output unchanged.
\end{theorem}

\begin{theorem}[Reset sweep turnaround]
\lean{Turing.ClockHost.p2_turn}
\label{Turing.ClockHost.p2_turn}
\leanok
\uses{Turing.ClockHost.apply_hostCfg, Turing.ClockHost.clockTr, Turing.ClockHost.hostCfg, Turing.ClockHost.hostCfg_symbols_yI, Turing.ClockHost.hostTM, Turing.ClockHost.stepWith_hostCfg, Turing.ClockHost.yard, Turing.ClockHost.yard_neg_one, Turing.NDTM.stepWith, Turing.moveInputPos}
\difficulty{3}
Let $S$ be a deterministic Turing machine over Boolean symbols, let $A$ be a
nondeterministic Turing machine over Boolean symbols, let $	ext{loop}$ be a Boolean, and
let $x$ be a Boolean input word. For every choice bit $b$, input-head position $ip$
among the $|x|+2$ positions, contents and head positions of the work tapes of both
machines, and natural number $m$, one step of the host machine from phase $p2$ with yard
tape storing $m$ consecutive true symbols and its head at $-1$ leaves $S$'s work tapes
and head positions unchanged, advances each of $A$'s work-tape heads by one without
changing their contents, and leaves the output empty. It moves the input head one
position to the right, clamped to the valid range, moves the yard head to $0$, and
enters phase $run$ with $A$ in its initial state and no captured output.
\end{theorem}

\begin{theorem}[Run during the first leftward reset sweep]
\lean{Turing.ClockHost.p0_marks}
\label{Turing.ClockHost.p0_marks}
\leanok
\uses{Turing.ClockHost.hostCfg, Turing.ClockHost.hostTM, Turing.ClockHost.inLeft, Turing.ClockHost.p0_mark, Turing.ClockHost.yard, Turing.ClockHost.yard_of_lt, Turing.NDTM.runWith, Turing.NDTM.runWith_cons}
\difficulty{4}
Let $S$ be a deterministic machine, $A$ a nondeterministic machine, and $x$ a Boolean
input word; let $m\in\mathbb N$, choose an input-head position and work-tape contents
and head positions for both machines, and let $w$ be a Boolean choice word of length at
most $m$. If the yard-tape head starts at $z=|w|-1$, then running the host machine from
phase $p0$ under $w$ leaves it in phase $p0$, with input-head position given by $|w|$
clamped left moves from its starting position, each of $A$’s work-tape heads advanced by
$|w|$, the yard-tape head at $-1$, and the other listed tape data and output unchanged.
\end{theorem}

\begin{theorem}[Repeated rightward sweep marks]
\lean{Turing.ClockHost.p1_marks}
\label{Turing.ClockHost.p1_marks}
\leanok
\uses{Turing.ClockHost.hostCfg, Turing.ClockHost.hostTM, Turing.ClockHost.p1_mark, Turing.ClockHost.yard, Turing.ClockHost.yard_of_lt, Turing.NDTM.runWith, Turing.NDTM.runWith_cons}
\difficulty{4}
Let $S$ be a deterministic Boolean-tape Turing machine and $A$ a nondeterministic
Boolean-tape Turing machine, with a Boolean loop parameter and input word $x$. For any
natural number $m$, choice word $w$ of length at most $m$, input-head position, and
simulated-machine tape contents and head positions, if the yard head starts at $m-|w|$,
then running the host machine from phase $p1$ under $w$ leaves it in phase $p1$,
advances each of $A$’s work-tape heads by $|w|$, places the yard head at $m$, and leaves
the other tapes and output unchanged.
\end{theorem}

\begin{theorem}[Repeated leftward sweep marks]
\lean{Turing.ClockHost.p2_marks}
\label{Turing.ClockHost.p2_marks}
\leanok
\uses{Turing.ClockHost.hostCfg, Turing.ClockHost.hostTM, Turing.ClockHost.p2_mark, Turing.ClockHost.yard, Turing.ClockHost.yard_of_lt, Turing.NDTM.runWith, Turing.NDTM.runWith_cons}
\difficulty{4}
Let $S$ be a deterministic Boolean-tape Turing machine and $A$ a nondeterministic
Boolean-tape Turing machine, with a Boolean loop parameter and input word $x$. For any
natural number $m$, choice word $w$ of length at most $m$, input-head position, and
simulated-machine tape contents and head positions, if the yard head starts at $|w|-1$,
then running the host machine from phase $p2$ under $w$ leaves it in phase $p2$,
advances each of $A$’s work-tape heads by $|w|$, places the yard head at $-1$, and
leaves the other tapes and output unchanged.
\end{theorem}

\begin{theorem}[Three-sweep reset run]
\lean{Turing.ClockHost.reset_run}
\label{Turing.ClockHost.reset_run}
\leanok
\uses{Turing.ClockHost.exists_split_block, Turing.ClockHost.hostCfg, Turing.ClockHost.hostTM, Turing.ClockHost.inLeft, Turing.ClockHost.inLeft_succ, Turing.ClockHost.p0_marks, Turing.ClockHost.p0_turn, Turing.ClockHost.p1_marks, Turing.ClockHost.p1_turn, Turing.ClockHost.p2_marks, Turing.ClockHost.p2_turn, Turing.ClockHost.yard, Turing.NDTM.runWith, Turing.NDTM.runWith_append, Turing.NDTM.runWith_cons, Turing.NDTM.runWith_nil, Turing.moveInputPos}
\difficulty{5}
Let $S$ be a deterministic Boolean-tape Turing machine and $A$ a nondeterministic
Boolean-tape Turing machine, with a Boolean loop parameter and input word $x$. For any
natural number $m$, input-head position, and simulated-machine tape contents and head
positions, if a choice word $w$ has length $3(m+1)$, then running the host machine under
$w$ from phase $p0$ with yard head at $m-1$ ends in the running phase with $A$ in its
initial state and no captured output; the input head is moved left by $m+1$ positions
with boundary clamping and then one position right, every head of $A$’s work tapes
advances by $3(m+1)$, the yard head is at $0$, and the other tapes and output are
unchanged.
\end{theorem}

\begin{definition}[Ready machine configuration]
\lean{Turing.ClockHost.ReadyAt}
\label{Turing.ClockHost.ReadyAt}
\leanok
\uses{Turing.Cfg}
Let $a$ be a configuration of the machine $A$ on input $x$, and let $m$ be a natural
number. The configuration is ready when its state is the initial state of $A$, its input
head is at position $1$, and its output is empty. In addition, on every work tape $i$,
every cell at position $z$ with $a$'s head position minus $m$ at most $z$ is blank.
\end{definition}

\begin{theorem}[Initial configuration is ready]
\lean{Turing.ClockHost.readyAt_initCfg}
\label{Turing.ClockHost.readyAt_initCfg}
\leanok
\uses{Turing.ClockHost.ReadyAt, Turing.NDTM.initCfg}
\difficulty{1}
Let $A$ be a nondeterministic Boolean-tape Turing machine, let $x$ be a Boolean input
word, and let $m$ be a natural number. The initial configuration of $A$ on $x$ is ready
at yard length $m$: it is in $A$’s initial state, its input head is at position $1$, its
output is empty, and every cell at or to the left of each work-tape head minus $m$ is
blank.
\end{theorem}

\begin{theorem}[State and Output Agreement for Short Runs]
\lean{Turing.ClockHost.ReadyAt.runWith_eq}
\label{Turing.ClockHost.ReadyAt.runWith_eq}
\leanok
\uses{Turing.Cfg, Turing.Cfg.init, Turing.ClockHost.ReadyAt, Turing.NDTM.AgreeShift, Turing.NDTM.AgreeShift.runWith, Turing.NDTM.initCfg, Turing.NDTM.runWith}
\difficulty{4}
Let $A$ be a nondeterministic finite-tape Turing machine, let $x$ be a Boolean input
word, and let $m$ be a nonnegative integer. If a configuration $a$ is ready at $m$ and a
Boolean choice word $u$ has length at most $m$, then running $A$ on $u$ from $a$ and
running it on $u$ from its initial configuration on $x$ yield configurations with the
same state and the same output word.
\end{theorem}

\begin{definition}[Next attempt configuration]
\lean{Turing.ClockHost.nextAttempt}
\label{Turing.ClockHost.nextAttempt}
\leanok
\uses{Turing.Cfg, Turing.ClockHost.inLeft, Turing.moveInputPos}
Let $a$ be a configuration of $A$ on input $x$, and let $m$ be a natural number. The
next-attempt configuration has the initial state of $A$, an input-head position obtained
by moving left $m+1$ steps from $a$'s input-head position and then right one step, the
same work-tape contents as $a$, and an empty output. Each work-tape head is shifted
right by $3(m+1)$ cells.
\end{definition}

\begin{theorem}[Readiness After a Choice-Word Run and Reset]
\lean{Turing.ClockHost.ReadyAt.next}
\label{Turing.ClockHost.ReadyAt.next}
\leanok
\uses{Turing.Cfg, Turing.ClockHost.ReadyAt, Turing.ClockHost.moveInputPos_inLeft, Turing.ClockHost.nextAttempt, Turing.NDTM.abs_runWith_workTapePos_sub_le, Turing.NDTM.runWith, Turing.NDTM.runWith_workTapes_of_le}
\difficulty{4}
Let $A$ be a nondeterministic finite-tape Turing machine, let $x$ be a Boolean input
word, and let $m$ be a nonnegative integer. If a configuration $a$ is ready at $m$, the
input length is at most $m$, and a Boolean choice word $u$ has length exactly $m$, then
the configuration obtained by running $A$ on $u$ from $a$ and applying the next-attempt
reset is ready at $m$.
\end{theorem}

\begin{theorem}[Work-head position bounds for a next attempt]
\lean{Turing.ClockHost.nextAttempt_workTapePos}
\label{Turing.ClockHost.nextAttempt_workTapePos}
\leanok
\uses{Turing.Cfg, Turing.ClockHost.nextAttempt, Turing.NDTM.abs_runWith_workTapePos_sub_le, Turing.NDTM.runWith}
\difficulty{3}
Let $A$ be a nondeterministic machine with $k$ work tapes, let $m\in\mathbb N$, let $a$
be a configuration of $A$, and let $u$ be a Boolean choice word of length $m$. For each
work-tape index $i$, the head position after the next-attempt reset applied to the run
from $a$ under $u$ lies between $a$’s initial position on tape $i$ plus $2m+3$ and that
position plus $4m+3$, inclusive.
\end{theorem}

\begin{definition}[Clock-ready configuration]
\lean{Turing.ClockHost.ClockReady}
\label{Turing.ClockHost.ClockReady}
\leanok
\uses{Turing.Cfg, Turing.ClockHost.ReadyAt, Turing.ClockHost.runCfg, Turing.ClockPhase}
Let $c$ be a configuration of the clock machine on input $x$, and let $m$ be a natural
number. The configuration $c$ is clock-ready exactly when there exist a configuration
$a$ of $A$ on input $x$, work-tape contents $s_T$ and work-head positions $s_P$ for the
yardstick machine, such that $a$ is ready and $c$ is the clock-machine configuration
obtained from $a$ at clock value $0$ with those yardstick tapes and head positions.
\end{definition}

\begin{theorem}[Nonhalting State of a Ready Clock Configuration]
\lean{Turing.ClockHost.ClockReady.state_ne_none}
\label{Turing.ClockHost.ClockReady.state_ne_none}
\leanok
\uses{Turing.Cfg, Turing.ClockHost.ClockReady, Turing.ClockHost.hostCfg, Turing.ClockHost.runCfg, Turing.ClockPhase}
\difficulty{3}
Let $S$ be a deterministic finite-tape Turing machine, let $A$ be a nondeterministic
finite-tape Turing machine, let $m$ be a nonnegative integer, and let $x$ be a Boolean
input word. Every clock-host configuration that is ready at $m$ has a nonempty state
component.
\end{theorem}

\begin{theorem}[Startup reaches the initial run configuration]
\lean{Turing.ClockHost.startup}
\label{Turing.ClockHost.startup}
\leanok
\uses{Turing.ClockHost.bootCfg, Turing.ClockHost.boot_run, Turing.ClockHost.hostCfg, Turing.ClockHost.hostTM, Turing.ClockHost.initCfg_eq_bootCfg, Turing.ClockHost.moveInputPos_inLeft, Turing.ClockHost.rewind_run, Turing.ClockHost.runCfg, Turing.ClockHost.yard, Turing.FinTM.ComputesInTime, Turing.FinTM.ComputesInTime.output_unique, Turing.FinTM.computesInTime_iff, Turing.MultiTapeTM.initCfg, Turing.MultiTapeTM.runFrom, Turing.NDTM.initCfg, Turing.NDTM.runWith, Turing.NDTM.runWith_append}
\difficulty{6}
Let $S$ be a finite-state deterministic Boolean machine, let $A$ be a nondeterministic
Boolean machine, let $x$ be a Boolean word, and let $m,T$ be natural numbers. Suppose
$S$ halts on input $x$ after exactly $T$ steps with output consisting of $m$ true
symbols, and $|x|\le m$. Then there is a natural number $s_0\le T+m+2$ and a choice of
work-tape contents and head positions for $S$ such that, for every Boolean word $w$ of
length $s_0$, running the host machine on input $x$ using choices $w$ for that many
steps yields the run configuration encoding the initial configuration of $A$ on $x$,
with clock value zero and those work-tape contents and head positions for $S$.
\end{theorem}

\begin{theorem}[Output preservation in one clocked attempt]
\lean{Turing.ClockHost.oneShot}
\label{Turing.ClockHost.oneShot}
\leanok
\uses{Turing.ClockHost.exists_split_block, Turing.ClockHost.hostCfg, Turing.ClockHost.hostTM, Turing.ClockHost.runCfg, Turing.ClockHost.run_run_zero, Turing.ClockHost.timeout_fail, Turing.ClockHost.timeout_success, Turing.NDTM.initCfg, Turing.NDTM.runWith, Turing.NDTM.runWith_append, Turing.NDTM.runWith_cons, Turing.NDTM.runWith_nil}
\difficulty{5}
Let $S$ be a deterministic machine and $A$ a nondeterministic machine over Boolean
symbols, let $x$ be a Boolean input word, and let $m\in\mathbb N$. For any work-tape
contents and head positions for $S$, and any Boolean choice word $w$ of length $m+1$,
the host machine with looping disabled, run from the corresponding run configuration
based on $A$’s initial configuration with the yard tape head at $0$, halts; moreover,
for every Boolean $v$, its output is exactly $[v]$ if and only if the run of $A$ under
the first $m$ bits of $w$ halts and has output exactly $[v]$.
\end{theorem}

\begin{theorem}[A Clocked Round]
\lean{Turing.ClockHost.round}
\label{Turing.ClockHost.round}
\leanok
\uses{Turing.Cfg, Turing.ClockHost.ClockReady, Turing.ClockHost.ReadyAt.next, Turing.ClockHost.ReadyAt.runWith_eq, Turing.ClockHost.exists_split_block, Turing.ClockHost.hostTM, Turing.ClockHost.nextAttempt, Turing.ClockHost.reset_run, Turing.ClockHost.run_run_zero, Turing.ClockHost.timeout_fail, Turing.ClockHost.timeout_success, Turing.ClockPhase, Turing.NDTM.initCfg, Turing.NDTM.runWith, Turing.NDTM.runWith_append, Turing.NDTM.runWith_cons, Turing.NDTM.runWith_of_halt}
\difficulty{4}
Let $S$ be a deterministic Boolean Turing machine, $A$ a nondeterministic Boolean Turing
machine, $x$ a Boolean input, and $m$ a natural number. Suppose a host configuration is
clock-ready, meaning it represents a ready configuration of $A$ with yard-head position
zero, and suppose $|x|\le m$. For every Boolean choice word $w$ of length $4(m+1)$, if
the run of $A$ on $x$ from its initial configuration under the first $m$ bits of $w$
halts with output consisting of the single bit $v$, then the host run halts with output
$v$; moreover, if that run of $A$ does not halt with a single-bit output, then the host
run is clock-ready.
\end{theorem}
